\section{Exact coefficients and comparison with ordinary graph Floer theory}
\label{cat:section}

The finite constructions of the preceding sections need not be compatible
as the energy cutoff increases.  We first pass their matrix coefficients
to a field in which all errors of unbounded valuation vanish.  We then
compare the resulting directed graph category with ordinary weakly exact
Floer theory.  This comparison concerns the graph objects only.  Its
purpose is to make the module constructed from \(L\) accessible to the
ordinary mirror theorem, without putting \(L\) into that theorem's
source category.

\subsection{The coefficient fields}

Write
\[
 \Lambda=
 \left\{\sum_{\lambda\in\R}a_\lambda\qnov^\lambda:
 a_\lambda\in\C,\quad
 \#\{\lambda<C:a_\lambda\ne0\}<\infty
 \text{ for every }C\in\R\right\}.
\]
The valuation of a nonzero series is the least exponent with nonzero
coefficient; set \(v(0)=+\infty\).  A nonempty support has a least element,
and the support condition makes addition and multiplication well defined.
After extracting its leading term, a nonzero series is invertible by a
geometric series.  Thus \(\Lambda\) is a field.

\begin{lemma}\label{cat:novikov}
The field \(\Lambda\) is complete and algebraically closed.  Its value group
is \(\R\), its residue field is \(\C\), and its cardinality is at most the
continuum.
\end{lemma}

\begin{proof}
In a Cauchy sequence, the coefficients below any fixed real bound
eventually stabilize.  Their stabilized values define a series satisfying
the same support condition, proving completeness.  The assertions about
the value group and residue field follow from the monomials and the
constant coefficient.

Here is a short algebraic-closure argument that also records why the
choice of complex coefficients matters.  Completeness implies that the
valuation is henselian.  Let \(E/\Lambda\) be a finite Galois extension of
degree \(d_E\).  Its uniquely extended valuation is Galois invariant.
For \(x\in E^\times\),
\[
 d_E\,v(x)=v\bigl(N_{E/\Lambda}(x)\bigr)\in\R.
\]
Divisibility of \(\R\), and torsion-freeness of an ordered group, imply
that \(v(x)\) already belongs to \(\R\).  The residue extension has degree
at most \(d_E\): linearly independent residue classes lift to linearly
independent elements after rescaling a putative relation to have a
coefficient of valuation zero.  It is therefore just \(\C\).
Consequently every Galois automorphism acts trivially on the residue
field.

For any \(0\ne x\in E\), choose \(a\in\Lambda^\times\) with \(v(a)=v(x)\).
All conjugates of \(x/a\) have the same nonzero residue.  Their sum has
nonzero residue, because the residue characteristic is zero.  Hence
\(\operatorname{Tr}_{E/\Lambda}(x)\ne0\).  The \(\Lambda\)-linear trace
therefore has zero kernel, which forces \(d_E=1\).  In characteristic
zero every finite extension lies in a finite Galois extension, so
\(\Lambda\) is algebraically closed.

Finally every support is countable, being a countable union of its finite
parts below integral bounds.  Both a countable list of real exponents and
a countable list of complex coefficients have at most continuum many
possibilities.
\end{proof}

Fix a nonprincipal ultrafilter \(\mathfrak u\) on \(\N\).  Let
\({}^*\Lambda=\Lambda^{\N}/\mathfrak u\), with its valuation in
\({}^*\R=\R^{\N}/\mathfrak u\).  Inequalities below mean inequalities on
an ultrafilter-large set.  Define
\begin{align}
 O_b&=\{x\in{}^*\Lambda:v(x)\ge-c
                         \text{ for some }c\in\R\},\nonumber\\
 \mathfrak m_\infty&=\{x\in{}^*\Lambda:v(x)>c
                         \text{ for every }c\in\R\},\nonumber\\
 \kappa&=O_b/\mathfrak m_\infty,\qquad R=\kappa[[t]].
 \label{cat:coefficient-rings}
\end{align}
Zero is included in both displayed subsets.

\begin{lemma}\label{cat:residue}
The ring \(O_b\) is a valuation ring with maximal ideal
\(\mathfrak m_\infty\).  Its residue field \(\kappa\) is algebraically
closed, and constant sequences induce an embedding
\(\Lambda\hookrightarrow\kappa\).  The fields and rings in
Equation~\eqref{cat:coefficient-rings} have cardinality at most the
continuum.

Suppose a fixed finite collection of polynomial identities in matrix
entries holds, for the \(N\)-th collection of entries, up to errors of
valuation greater than \(N\).  If all those entry sequences are bounded
below in valuation, their images in \(\kappa\) satisfy the identities
exactly.  The same assertion holds coefficientwise for formal power
series in \(t\).
\end{lemma}

\begin{proof}
The elements of \({}^*\R\) bounded above and below by standard reals form
a convex subgroup.  Coarsening the valuation by this subgroup gives
exactly \(O_b\); positive coarsened valuation gives
\(\mathfrak m_\infty\).  Alternatively, an element of \(O_b\) outside
\(\mathfrak m_\infty\) has valuation bounded on both sides, so its inverse
also belongs to \(O_b\).  Every nonzero constant sequence has this
property, proving the asserted embedding.

The ultrapower is algebraically closed: a polynomial has finitely many
coefficients, and choosing a root coordinatewise proves this assertion
directly.  A monic polynomial over \(\kappa\) lifts to a monic polynomial
over \(O_b\).  Each root of the lift in \({}^*\Lambda\) lies in \(O_b\);
otherwise its highest power would have strictly smaller coarsened
valuation than every other term.  Reducing the factorization proves
algebraic closure of \(\kappa\).  The cardinality bound follows from
\((2^{\aleph_0})^{\aleph_0}=2^{\aleph_0}\).

Reduction \(O_b\to\kappa\) respects finite sums and products, and a sequence
with valuation greater than \(N\) belongs to \(\mathfrak m_\infty\).
For a fixed coefficient of \(t\), a polynomial expression in power
series uses only finitely many coefficients.  This proves the last
assertion.
\end{proof}

\subsection{The exact directed category and its module}

Choose finite lists exhausting the graph labels, and let the energy and
external-arity cutoffs tend to infinity through \(N\).
The almost complex data, cap lines, and generator conventions are fixed
once and for all; only the finite perturbations and the solutions
\(b_{L,N}\) of the curvature equation are chosen anew.
Apply Corollary~\ref{obs:finite-module}, and truncate each post-twist
matrix entry through energy \(N\).  For any fixed word, all sufficiently
large lists contain that word.  Entries for the finitely many missing
lists may be set equal to zero.

\begin{proposition}\label{cat:module}
These choices produce the following uncurved, nonunital
\(\Ainf\) data:
\begin{enumerate}[label=\textup{(\roman*)}]
 \item an integer-graded directed category \(\mathcal D_\Lambda\), with
 objects all graph labels and
 \[
 \Hom_{\mathcal D_\Lambda}(X,Y)=
 \begin{cases}
   \displaystyle\bigoplus_{x\in X\cap Y}\Lambda x,&X\prec Y,\\
   0,&\text{otherwise};
 \end{cases}
 \]
 \item its scalar extension \(\mathcal D_\kappa\), and a right
 \(\Z/4\)-graded module \(\mathcal M_0\) over the reduction modulo four
 of \(\mathcal D_\kappa\), with value complex
 \[
 \mathcal M_0(X)=\bigoplus_{x\in X\cap L}\kappa x;
 \]
 \item on the labels with scalar slope, a weighted integer-graded
 directed category \(\mathcal D_R\), and a right \(\Z/4\)-graded module
 \(\mathcal M_R\), with finite free value complexes
 \[
 \mathcal M_R(X)=\bigoplus_{x\in X\cap L}R x.
 \]
 Reduction modulo \(t\) identifies \(\mathcal D_R\) and \(\mathcal M_R\)
 with the corresponding restrictions of \(\mathcal D_\kappa\) and
 \(\mathcal M_0\), including all operations.
\end{enumerate}
The graph category is given by actual Novikov-convergent ordinary signed
counts.  In particular it does not depend on any asserted convergence
of the sequence \(b_{L,N}\).
\end{proposition}

\begin{proof}
All Taylor coefficients of the finite post-twist entries have
nonnegative valuation.  Indeed unforced components have nonnegative
area, and each inserted \(b_{L,N}\) has valuation at least
\(\delta_L>0\).  For each fixed arity, the category and module equations
are finite matrix identities.  Corollary~\ref{obs:finite-module}
makes their errors have valuation greater than \(N\), coefficientwise
in \(t\), for all sufficiently large \(N\).
Lemma~\ref{cat:residue} therefore makes them exact after reduction.
Taking the constant Taylor coefficient commutes with these polynomial
identities.  The common choice of the finite weighted and unweighted
systems gives the exact specialization assertion.  No comparison
between \(b_{L,N}\) and \(b_{L,N+1}\) is used.

For a directed word, all rigid moduli are regular, and
Theorem~\ref{virt:finite-integration} identifies the finite virtual
integral with their ordinary signed count.  Compactness gives only
finitely many contributing classes below a fixed energy bound.
Consequently the count defines an element of \(\Lambda\), or
coefficientwise of \(\Lambda[[t]]\) when the relative weight is present.
Its difference from the finite-cut entry has valuation tending to
infinity.  Thus its image in \(\kappa\), respectively \(R\), is precisely
the coefficient just constructed.  Integer degrees on the graph
category and modulo-four degrees on its module follow from
Lemma~\ref{obs:degree}.
\end{proof}

We write the module action in boundary order.  An arrow \(x:X\to Y\)
acts contravariantly on values:
\[
 \mathcal M_0(Y)\longrightarrow\mathcal M_0(X).
\]
Our suspended convention has insertion sign
\(\sum_{j<i}(|x_j|-1)\), as in Proposition~\ref{obs:ainfinity}.
On cohomology the associative, unsuspended composition of successive
arrows \(x,y\) is \((-1)^{|x|}m_2(x,y)\).
In particular no additional sign occurs when the first arrow has degree
zero.  The same convention applies to the binary module action.

\begin{lemma}\label{cat:euler}
For every graph label \(X\),
\begin{equation}\label{cat:euler-equation}
 \sum_{r\in\Z/4}(-1)^r
       \dim_\kappa H^r(\mathcal M_0(X))
   =[X]\cdot[L]=\nu[X]\cdot\alpha.
\end{equation}
For ordinary integer-graded graph Floer complexes the analogous Euler
characteristic is the oriented intersection number.
\end{lemma}

\begin{proof}
A finite-dimensional complex with an even grading period has the same
alternating dimension as its cohomology: each boundary space appears
twice with opposite signs.  It remains to check the sign of an
intersection generator.

At a transverse intersection choose unitary coordinates with the first
tangent plane real.  Let \(\alpha_1,\ldots,\alpha_8\in(0,\pi)\) be the
positive plane-rotation angles, and let \(\vartheta_1,\vartheta_2\) be
the grading lifts, taken modulo four when necessary.  The generator
degree is
\begin{equation}\label{cat:cap-degree}
 |x|=8+\vartheta_2-\vartheta_1
                  -\sum_{\ell=1}^8\alpha_\ell/\pi.
\end{equation}
The positive-rotation basis of the second plane differs from its marked
orientation by
\((-1)^{\vartheta_2-\vartheta_1-\sum\alpha_\ell/\pi}\).
Concatenating the first-plane basis with that rotation basis has ambient
symplectic orientation sign \((-1)^{8\cdot7/2}=1\).  Since eight is even,
the local intersection sign is exactly \((-1)^{|x|}\).
Summing proves the result, using \([L]=\nu\alpha\).
\end{proof}

\subsection{Ordinary graph data extending the prescribed sector}

An ordinary category used below has all ordered words of graph objects,
including repeated labels; its morphisms for repeated or nontransverse
positions use Hamiltonian chords.  The relevant established ingredients
are the transverse Floer compactness, determinant orientations,
continuation, and gluing constructions; see, for example,
\cite{seidel}.  We must verify that our prescribed domain-dependent
data can be included without changing their strict-word counts.
The verification also explains the closed weakly exact setting used
here.

\begin{proposition}\label{cat:ordinary-comparison}
There is an ordinary integer-graded, zero-background spin, weakly exact
Floer category \(\mathcal C\) of the graph labels and a nonunital
\(\Ainf\) functor
\[
 \mathcal D_\Lambda\longrightarrow\mathcal C
\]
which is the identity on objects and induces an isomorphism on morphism
cohomology whenever \(X\prec Y\).
Every finite collection in \(\mathcal C\) is comparable, by ordinary
quasi-equivalences, with the corresponding ordinary graph category
constructed from any other regular compatible Floer data.
These comparisons require no polygons with boundary on \(L\).
\end{proposition}

\begin{proof}
We first extend the data, then establish compactness and gluing, and
finally compare conventions and choices.

\paragraph{Data on the compactified domain families.}
For each additional ordered pair choose regular strip data
\((H_\tau\,d\tau,J_\tau)\) with nondegenerate chords.  On strict pairs
retain exactly the unforced pair data already chosen.
For a word of length at least three, use a Hamiltonian-valued one-form
\(\mathsf K\) and a compatible \(J\) on its punctured disks.
Require the tangential restriction of \(\mathsf K\) on each boundary arc
to vanish.  On every strict word use the original \(\mathsf K=0\)
harmonic data, unchanged.

The requisite extension space can be described explicitly.
Replace each puncture by its strip-height interval.  Near a boundary
node use the real plumbing coordinates
\[
 r_l=e^{-\pi s},\qquad r_r=e^{-\pi(T-s)},\qquad
 r_l r_r=\rho=e^{-\pi T},\qquad 0\le\tau\le1.
\]
We require smooth coefficients in
\((r_l,r_r,\tau)\), with the other component parameters as spectators.
For one-forms, smoothness refers to their coefficients in the vertical
strip coframe, dual to
\[
 -r_l\partial_{r_l}+r_r\partial_{r_r},
 \qquad \partial_\tau.
\]
At the added interval prescribe the pair data
\((H_\tau\,d\tau,J_\tau)\); on the two nodal faces prescribe the
component data.  At external ends use the one-sided version.

These conditions are independent of the chosen plumbing chart.
Indeed a change of holomorphic boundary-node coordinate has positive
real linear term.  Its radial variables therefore change by smooth
positive factors.  A displacement of a node origin in a smoothed fiber
is \(O(\rho)\); dividing by \(r_l\) produces \(O(r_r)\), and conversely.
Thus the vertical strip transitions extend smoothly and have the same
limiting height and strip directions.  This is also a direct description
of the real radial blowup of the boundary plumbing.

The harmonic data obey these conditions by
Lemma~\ref{local:neck-coefficients}.
For any new word, the prescribed face restrictions agree at every
intersection: both restrictions are the data attached to the same
component tree.  They can consequently be extended over corner charts
by extension of face functions and inclusion--exclusion, and then
patched by a partition of unity.  The zero tangential coefficient on a
boundary arc is preserved by extending that coefficient with its
boundary defining factor.  Compatible structures may be interpolated
in a local chart, or by interpolating metrics followed by compatible
retraction.  This inductive extension never changes a strict word,
since every component of a strict word is again strict.

On new stable words, Hamiltonian one-form variations supported at an
interior domain point yield the usual universal transversality.
Their Hamiltonian vector fields span the target tangent space, and
domain dependence localizes the variation regardless of whether the
map is multiply covered.  Elliptic unique continuation then gives
surjectivity of the universal linearization.  Relative Sard--Smale
choices give regularity while preserving prescribed faces.  There are
only countably many words and relative classes, so the choices can
be made simultaneously.  Strict words already have the required
regularity.

\paragraph{Energy and compactness.}
On a neck, the difference from the pair model is
\(O(r_l+r_r)\), with all derivatives needed here.  The curvature of
\(H_\tau\,d\tau\) is zero.  Thus the Hamiltonian curvature of
\(\mathsf K\) has an integral supremum bound uniform in the conformal
parameters of a fixed word:
\[
 \int_R^{T-R}\bigl(e^{-\pi s}+e^{-\pi(T-s)}\bigr)\,ds
 \le \frac{2}{\pi}e^{-\pi R}.
\]
The finitely many core charts give the remaining bound.  The energy
identity for \(du-X_{\mathsf K}\) therefore differs from
\(\int u^*\omega\) by a uniformly bounded curvature term and the fixed
external chord terms.  In particular the topological area is bounded
below, and an upper area bound gives a uniform analytic energy bound.

There are no nonconstant unforced spheres, since \(\pi_2(T)=0\).
There are no unforced disks on an individual graph, since
\(\pi_1(X)\to\pi_1(T)\) is injective and
\(\pi_2(T,X)=0\).  Small-scale bubbles of an inhomogeneous equation would
be exactly such unforced disks or spheres, so none occur.
Translation on a growing neck limits to the prescribed pair strip
equation.  Nonconstant strip breakings have a positive energy gap for
each fixed pair.  On the remaining small-energy neck portions,
solutions approach isolated chords; Hamiltonian transport of the pair
equation reduces the estimate to a transverse intersection with
height-dependent almost complex structure.  Its asymptotic operator
has a spectral gap.  The exponential estimates then apply, with the
exponentially decaying \(O(r_l+r_r)\) coefficient errors included.
Thus necks lose neither energy nor matching data.

The spectral-gap and neck estimates are those of
Lemma~\ref{local:decay}, after Hamiltonian transport.
This is the ordinary compactness argument, with precisely the same
domain degenerations as the compactified word family.  The areas are
homotopy invariant for the fixed chord ends and add on broken trees:
the matched chord intervals occur with opposite orientations.  Core
closeness and the small neck neighborhoods also give eventual homotopy
to the limiting tree, so only finitely many relative classes occur
below a fixed area bound.  Consequently the rigid counts define
Novikov series and the one-dimensional compactifications have exactly
the expected broken boundaries.

For later use one may allow a nonzero tangential component on boundary
arcs, provided it has the same prescribed pair limits and face
restrictions.  Its coefficient in the strip-length direction is then
\(O(r_l+r_r)\) near the added intervals; its integral, and hence the
additional boundary term in the energy identity, remains uniformly
bounded.  The preceding conclusions still apply.

\paragraph{Gluing.}
The transverse gluing construction applies to these data as well.
On the left half of a length-\(T\) neck the difference from the
partially broken left datum is \(O(r_r)\), and on the right half it is
\(O(r_l)\).  In coordinates transporting the pair Hamiltonian flow,
the limiting equation is the transverse unforced strip equation;
the additional forcing at its constant solution is
\(O(r_l+r_r)\).  Choose an exponential weight smaller than both the
asymptotic spectral gap and the coefficient decay rate.  Cutting off
the decaying component solutions gives an exponentially small
pregluing error in this weight.

To construct a right inverse, split a forcing term in a fixed-width
region around the middle of each neck, solve on the two component
ends, and splice the resulting sections with cutoffs whose derivatives
are \(O(1/T)\).  The component weights dominate the joined weight.
The error consists of these cutoff derivatives and the coefficient
differences just estimated; its operator norm tends to zero.
A Neumann series therefore corrects the parametrix.  Fixing core kernel
projections makes the augmented linearization invertible, with a
uniform bound.  The ordinary uniform nonlinear estimates then give
the unique nearby glued solution.  Exponential decay gives the
converse: every nearby curve lies in this gluing neighborhood.
The same construction can leave any subset of nodes broken.

Length derivatives in coordinates fixed on the corresponding
component halves have exponential bounds: explicit coefficient
derivatives contain the opposite-end radial factor, while
reparametrization derivatives in the middle multiply exponentially
decaying map derivatives.  Differentiating the inverse construction
in slightly weaker weights gives the boundary smoothness needed for
zero- and one-dimensional moduli.  Thus the compactification and
orientation gluing give the ordinary \(\Ainf\) relations.  This is the
transverse case of the splice used earlier, with no movable
smooth-node mode.  It requires no virtual perturbation for the graph
category.

\paragraph{Comparison of choices.}
Two systems can be put on two differently labeled copies of their
object lists.  Prescribe each system on its own words, and extend on
mixed words by the preceding construction.  There is no conflict on
faces.  Ordinary data with exact strip collars are admissible here too.
Interpolating strip data with fixed external ends define continuation
maps.  Their homotopies have bounded curvature on compact interpolation
regions and the same no-bubbling property.  Gluing interpolations and
deforming their concatenation proves composition and independence.
Attaching continuation strips at the three triangle ends, and varying
the remaining compact data, proves compatibility with binary products.

It follows that the inclusions of either copy are cohomologically full
and faithful, and the two copies of each object are isomorphic on
cohomology.  The resulting quasi-equivalences are the ordinary
comparison of data.  Ordinary self-Floer units give cohomological
units in the combined category.  The same argument includes
Hamiltonian-transported copies, as made explicit in
Lemma~\ref{cat:gauge} below.

For the strict-word inclusion one finally identifies determinant
conventions.  Lemma~\ref{cat:sign-comparison} proves that this is a
strict change of generator signs.  The equations on the strict
sector are then literally those of \(\mathcal D_\Lambda\).
This completes the construction and comparison.
\end{proof}

The proof applies as well after adjoining finitely many of the auxiliary
vertical tori used in Lemma~\ref{local:comparison}.  They are graded spin
tori with injective fundamental-group maps, so the same disk and sphere
exclusions, regular perturbations, energy estimates, and continuation
argument apply.  We use this extension only to compute the comparison
triangle products below.

\subsection{Hamiltonian gauge and conventions}

\begin{lemma}\label{cat:gauge}
Replacing any finite set of graph labels by specified Hamiltonian
images, with their transported gradings and spin structures, preserves
the ordinary category up to the above comparisons.  Under the
corresponding identification of polygons, topological area weights
differ only by fixed input and output factors
\(\qnov^{a_x}\), one for each chord.
\end{lemma}

\begin{proof}
On a word disk let \(G_z(v,x)\), \(0\le v\le1\), be normalized
Hamiltonians, equal on each open boundary arc to the Hamiltonian
specified for that label.  The functions can be interpolated using
harmonic arc weights indexed by arc positions; the face and
radial-collar rules are the same as above.  Let \(\phi_z^v\) be their
time-\(v\) flows and put
\[
 F_v(z,x)=(z,\phi_z^v(x)),\qquad
 \Gamma=\int_0^1 F_v^*(d_zG_z(v,\cdot))\,dv .
\]
Here \(d_z\) differentiates with the target variable held fixed.
With the convention \(\iota_{X_G}\omega=-d_xG\), Cartan's formula gives
\begin{equation}\label{cat:gauge-equation}
 F_1^*\omega=\omega+d\Gamma.
\end{equation}
Indeed
\[
 \frac d{dv}F_v^*\omega
 =dF_v^*(-d_xG_z)
 =dF_v^*(d_zG_z),
\]
since the omitted full differential has zero exterior differential.

Conjugating maps by \(\phi_z^1\) gives Floer data of the admitted kind.
Every variation of a Hamiltonian flow is Hamiltonian.  The moduli
spaces therefore correspond, with their determinant orientations
transported using the isotopy of the generator caps.  On an open
boundary arc the generating Hamiltonian is independent of its arc
coordinate, so \(\Gamma\) vanishes on the boundary tangent.
Stokes in Equation~\eqref{cat:gauge-equation} leaves only the integrals
over external chord intervals.  Those depend just on the fixed chord,
with the input and output signs dictated by boundary orientation.
They are absorbed by rescaling the chord basis.

One obtains the same conclusion when ordinary perturbations are
presented by moving boundaries.  Extend their generating functions over
the disk with the pair limits and the face conditions, and conjugate
back to fixed boundaries.  The resulting one-form may have a tangential
component, which is allowed by the energy estimate above.  On a
boundary tangent \(V\), \(\Gamma(V)\) is the normalized Hamiltonian
generating \(\partial_V\phi_z^1\), pulled back by \(\phi_z^1\).
This follows either from the mixed tangent component of
Equation~\eqref{cat:gauge-equation} or directly by differentiating the
flow.  The moving-boundary Hamiltonian correction therefore gives the
same fixed-boundary area up to the stated end factors.
\end{proof}

There is a harmless variance choice in comparing Floer conventions.
Our strip has the previous label at height zero and the next at height
one, with inputs at positive ends.  Its input degree is the index of
the closing \emph{negative}-end cap, giving
Equation~\eqref{cat:cap-degree}.  Changing tensor notation does not
change that index.  If a reference orders the geometric boundary
labels oppositely, use its \(\Ainf\) opposite, reversing tensors with
the suspension Koszul signs.  These are the same geometric disks;
no reflection of a holomorphic map is involved.  For example, with
matched grading lifts, the degree assigned to the opposite corner is
\[
 \vartheta_A-\vartheta_B
       +\sum_\ell\alpha_{A,B,\ell}/\pi
 =8+\vartheta_A-\vartheta_B
       -\sum_\ell\alpha_{B,A,\ell}/\pi.
\]
Thus it agrees with our cap index for the arrow \(B\to A\).
On the mirror side a variance reversal is matched by perfect duality.
The ordinary graph theorem may consequently be used in the
covariant form needed below.

\begin{lemma}\label{cat:sign-comparison}
After matching cap determinant lines, the strict directed operations
of Section~\ref{cat:section} and ordinary spin Floer operations with
the same boundary and suspension order are strictly isomorphic by
sign changes of generators.  These changes preserve degrees and
Novikov valuations.
\end{lemma}

\begin{proof}
Use the same negative- and positive-end cappings, reciprocal
contractions, and spin trivializations for the two constructions.
A change of basis at one end is tracked also in the cap contraction;
the remaining freedom is precisely the ordinary sign of the generator
line.  Having matched those lines, write
\[
 (-1)^{r_k(a_1,\ldots,a_k)}
\]
for the ratio of the two signs on a rigid \(k\)-input operation, where
\(a_j=|x_j|-1\) is the suspended integer degree.
The output suspended degree is \(1+\sum_j a_j\).

We justify that this ratio depends only on these indices.
Both constructions orient the capped disk by its spin boundary
framing and orient sphere factors complexly.  There is no background
class.  Integral graph gradings make the capped Maslov contribution
zero.  The two rules use the same spin homotopy class along each
boundary segment, so their ratio acquires no relative-class twist.
At fixed degree, a transverse graded pair can be varied through
connected data: after fixing the first unitary frame, deform the
second transverse plane to the imaginary plane in its graph chart,
transporting its lift and orientation.  Positive changes of frame
are connected.  Along a segment, lifted determinant winding determines
the homotopy class of a unitary frame path, since \(\SU(8)\) is simply
connected.  Capping and lower-order changes therefore leave only
permutations of determinant, marking, and suspension factors; those
depend solely on the indices and the order of the inputs.
The same argument applies to the one-dimensional determinant
comparisons.

Both conventions have the same suspended insertion sign.
Determinant gluing at a two-vertex boundary consequently implies that
\begin{equation}\label{cat:sign-cocycle}
 \begin{split}
 &r_s(a_i,\ldots,a_{i+s-1})\\
 &\quad+
 r_{k-s+1}\left(a_1,\ldots,a_{i-1},
         1+\sum_{j=i}^{i+s-1}a_j,
         a_{i+s},\ldots,a_k\right)
 \end{split}
\end{equation}
is independent of the chosen block and its length \(s\ge1\).
This is a determinant-line identity, not an inference from a
potentially empty moduli space.  It holds for arbitrary integer
tuples: choose the ordered graded end planes with the prescribed
indices, choose the node and output degrees for two rigid formal
vertices, and join the boundary paths in the graded-plane cover.
The one-dimensional gluing comparison then gives
Equation~\eqref{cat:sign-cocycle}; strip vertices include their
translation direction in this comparison.

Choose \(f:\Z\to\Z/2\) recursively so that
\[
 f(a+1)+f(a)=r_1(a).
\]
Rescaling a generator of suspended degree \(a\) by \((-1)^{f(a)}\)
replaces \(r_k\) by
\[
 r'_k=r_k+f(1+\textstyle\sum_j a_j)+\sum_j f(a_j).
\]
It makes \(r'_1=0\) and preserves the block identity.
Compare a unary child at input \(i\) with a unary parent.  The identity
then gives
\[
 r'_k(a_1,\ldots,a_i+1,\ldots,a_k)
     =r'_k(a_1,\ldots,a_i,\ldots,a_k).
\]
Hence \(r'_k\) is a constant \(c_k\).  Comparing a binary child with a
unary parent gives \(c_k=c_2+c_{k-1}\), with \(c_1=0\); thus
\(c_k=(k-1)c_2\).  A common sign change of every arrow by
\((-1)^{c_2}\) changes a \(k\)-input operation by
\((-1)^{(k+1)c_2}=(-1)^{(k-1)c_2}\), removing all remaining signs.
The resulting identification is strict.
\end{proof}

\subsection{Later copies act invertibly}

For a graph label \(X\), let \(X^+\) be its next-layer copy.
The close perfect Morse perturbation has a unique degree-zero
generator, the maximum, denoted
\[
 w_X\in\Hom_{\mathcal D_\Lambda}^0(X,X^+).
\]
This is the arrow used later to pass from the directed category to
the full line-bundle model.

\begin{proposition}\label{cat:continuation}
The arrow \(w_X\) is closed.  In the ordinary graph category its class is
an isomorphism, and its action is a chain isomorphism
\[
 \mathcal M_0(X^+)\xrightarrow{\;w_X\;}\mathcal M_0(X).
\]
In particular it induces an isomorphism in each modulo-four cohomology
degree.
\end{proposition}

\begin{proof}
Self-Floer cohomology of a weakly exact graph with the transported brane
structure is its ordinary cohomology over \(\Lambda\).  For a sufficiently
small exact perturbation this is the usual Morse comparison; ordinary
continuation identifies all the choices, and no disk or sphere terms
occur.  The chosen perfect Morse pair has exactly the total Betti
number of generators.  Its differential must therefore be zero:
the total dimension of cohomology equals that of the underlying finite
complex.  The degree-zero group is one-dimensional.  Its nonzero
generator \(w_X\) is consequently a nonzero scalar multiple of a
Hamiltonian continuation isomorphism, and hence itself an isomorphism.
The sign identification above preserves these assertions.

We now prove the stronger statement about the newly constructed
module.  Put \(Y=X^+\), and use Lemma~\ref{local:comparison}:
\(X\) is horizontal, and near every point of
\(X\cap L\) the other sheet is a vertical patch \(q=q_j\).
The graph \(Y\) is \(p=\nabla f(q)\), with oscillation less than
\(\epsilon_X\).  Here \(\epsilon_X\) is smaller than both the disk gap
\(\delta_L\) and the fixed crossing energy needed to leave the chosen
tube.  The intersections of \(X\) and \(Y\) with \(L\) are paired by
these vertical patches.  The generator \(w_X\) is at the unique maximum
of \(f\), with the degree and brane structures transported through the
coordinate change.

Let \(F_j^{\mathrm{vert}}\) be the full comparison vertical torus at
\(q_j\).  Every triangle for
\((X,Y,F_j^{\mathrm{vert}})\) with first input \(w_X\) has area
\begin{equation}\label{cat:diagonal-area}
 \Delta_j=f(\max)-f(q_j)\le\operatorname{osc}(f)<\epsilon_X,
\end{equation}
by lifting to the cover and applying Stokes.
The two Floer complexes with the vertical torus each have one
generator and no differential.  Product with the isomorphism \(w_X\)
is therefore nonzero.  Its coefficient is
\[
 c_j\qnov^{\Delta_j},\qquad c_j\in\C^\times.
\]
All triangles contributing to it have the single area in
Equation~\eqref{cat:diagonal-area}, so this assertion concerns the
ordinary signed count itself.

The data of Proposition~\ref{local:data} for these comparison triangles
were chosen to agree with
those for \((X,Y,L)\) inside the tube.  Every such low-area triangle
stays in one vertical patch; there are none between unmatched patches.
There are also no nonconstant low-area strips for \((X,L)\) or
\((Y,L)\): inside the tube their ends belong to the same patch and
Stokes gives area zero.  Thus replacing the comparison torus by \(L\)
identifies the low-area rigid triangle spaces literally.
Their orientations agree up to a fixed nonzero sign at each patch.
Indeed the boundary arc lies in a contractible common vertical patch,
so its spin trivialization and caps can be identified, allowing only
a fixed grading shift and basis sign, independent of the triangle.
Theorem~\ref{virt:finite-integration} makes the finite virtual count in
this regular sector equal to the ordinary signed count.

No term with a bounding-cochain insertion occurs at these energies,
because each insertion has valuation at least
\(\delta_L>\epsilon_X\).
Consequently, for every sufficiently large cutoff, the matrix \(A_N\)
of action by \(w_X\) has the form
\[
 A_N=\diag(c'_j\qnov^{\Delta_j})+E_N,
 \qquad c'_j=\pm c_j\ne0,\qquad
 v((E_N)_{ij})>\epsilon_X.
\]
The determinant's uniquely lowest-valuation term is the product of the
displayed diagonal entries.  In fact any changed row of a determinant
term replaces a factor of valuation \(\Delta_j<\epsilon_X\) by one of
valuation greater than \(\epsilon_X\).  Diagonal corrections have the
same strict inequality.  This comparison is row by row; it does not
require \(\sum_j\Delta_j<\epsilon_X\).
Thus
\[
 v(\det A_N)=\sum_j\Delta_j
\]
for every sufficiently large \(N\), with the fixed nonzero leading
coefficient \(\prod_j c'_j\).  In particular \(\det A_N\) is not in
\(\mathfrak m_\infty\), so the limiting matrix over \(\kappa\) is
invertible.  The exact module identity makes it a degree-zero chain
map, hence a chain isomorphism.

If \(X\cap L\) is empty, \(Y\) was chosen disjoint from \(L\) as well.
Both value complexes are then zero, and the assertion holds.
\end{proof}
